%------------------------------------------------------------------------------%
% Luis A. D'Afonseca
%
%------------------------------------------------------------------------------%
\documentclass[a4paper,12pt,fleqn]{article}

\usepackage{style_avaliacao}

\ShowAnswers

%------------------------------------------------------------------------------%
\begin{document}

\makeHeader{Integrais e Séries}{Prova 4 -- Turma 5}{12/12/2023}

% 01
%------------------------------------------------------------------------------%
\exercise{30}
Calcule a Série de Taylor da função
\(
  f(x) = \int e^{-x^2} dx
\)

\begin{answer}
  Sabemos que a Série de Taylor da função exponencial é
  \[
    e^x = \sum_{n=0}^\infty \frac{x^n}{n!}
  \]
  Primeiro calculamos a Série de Taylor da função
  \[
    e^{-x^2}
    = \sum_{n=0}^\infty \frac{\left(-x^2\right)^n}{n!}
    = \sum_{n=0}^\infty \frac{(-1)^n}{n!} x^{2n}
  \]
  Agora integramos
  \begin{align*}
    f(x)
    & = \int e^{-x^2} dx \\[2mm]
    & = \int \sum_{n=0}^\infty \frac{(-1)^n}{n!} x^{2n} \;dx \\[2mm]
    & = \sum_{n=0}^\infty \frac{(-1)^n}{n!} \int x^{2n} dx \\[2mm]
    & = C + \sum_{n=0}^\infty \frac{(-1)^n}{n!} \frac{x^{2n+1}}{2n+1} \\[2mm]
    & = C + \sum_{n=0}^\infty \frac{(-1)^nx^{2n+1}}{(2n+1)n!}
  \end{align*}
\end{answer}

% 02
%------------------------------------------------------------------------------%
\exercise{40}
Seja \( f(x) = \ln(x+1)\)
\begin{enumerate}
  \item[{a) [12]}] Encontre um padrão para \(f^{(n)}(x)\)
  \item[{b) [14]}] Usando o padrão do item anterior, encontre a
                   Serie de Taylor de $f(x)$, centrada em zero
  \item[{c) [14]}] Obtenha o raio de convergência da série de potências
\end{enumerate}

\begin{answer}
  \begin{enumerate}[label=\alph*)]
    \item Calculando as primeiras derivadas
    \begin{align*}
      f^{(0)}(x) & = \ln(x+1) \\[2mm]
      f^{(1)}(x) & = \frac{\,d\;}{\,dx\;}\ln(x+1) = \frac{1}{x+1}(x+1)' = (x+1)^{-1} \\[2mm]
      f^{(2)}(x) & = \frac{d^2}{dx^2}    \ln(x+1) =    -                  (x+1)^{-2} \\[2mm]
      f^{(3)}(x) & = \frac{d^3}{dx^3}    \ln(x+1) = \hpm 2              \,(x+1)^{-3} \\[2mm]
      f^{(4)}(x) & = \frac{d^4}{dx^4}    \ln(x+1) =    - 2\cdot 3       \,(x+1)^{-4} \\[2mm]
      f^{(5)}(x) & = \frac{d^5}{dx^5}    \ln(x+1) = \hpm 2\cdot 3\cdot 4\,(x+1)^{-5}
    \end{align*}
    Vemos que, para $\bm{n\geq1}$,
    \[
      f^{(n)}(x)
      = \frac{d^n}{dx^n} \ln(x+1)
      = (-1)^{n+1} (n-1)! (x+1)^{-n}
      = \frac{(-1)^{n+1} (n-1)!}{(x+1)^n}
    \]
    \item O primeiro coeficiente, $n=0$, da Série de Taylor é
      \[
        c_0 = \frac{\ln(0+1)}{0!} = \frac{0}{1} = 0
      \]
      portanto, a série começa com $n=1$.
      Os demais coeficientes, $n\geq 1$, são
      \[
        c_n
        = \frac{f^{(n)}(0)}{n!}
        = \frac{(-1)^{n+1} (n-1)!}{(0+1)^nn!}
        = \frac{(-1)^{n+1}}{n}
      \]
      Montando a série temos
      \[
        T(x)
        = \sum_{n=0}^\infty c_n x^n
        = \sum_{n=1}^\infty \frac{(-1)^{n+1}}{n} x^n
      \]
      \item Aplicando o Teste da Razão temos
      \begin{align*}
        \rho
        & = \lim_{n\to\infty} \frac{\abs{u_{n+1}}}{\abs{u_n}} \\[2mm]
        & = \lim_{n\to\infty} \abs{\frac{(-1)^{n+2} x^{n+1}}{n+1}} \abs{\frac{n}{(-1)^{n+1} x^n}} \\[2mm]
        & = \lim_{n\to\infty} \frac{\abs{x}^{n+1}}{n+1} \frac{n}{\abs{x}^n} \\[2mm]
        & = \lim_{n\to\infty} \frac{n}{n+1} \abs{x}\\[2mm]
        & = \abs{x} \lim_{n\to\infty} \frac{n}{n+1}\\[2mm]
        & = \abs{x}
      \end{align*}
      Impondo a condição de convergência $\rho<1$ temos
      \[
        \rho = \abs{x} < 1
      \]
      portanto o raio de convergência é $R=1$


  \end{enumerate}
\end{answer}

% 03
%------------------------------------------------------------------------------%
\exercise{30}
Prove que a série de Maclaurin de $e^x$
converge para a própria função, para todo $x\in\R$

\begin{answer}
  Pelo \textbf{Teorema de Taylor}, sabemos que o erro cometido ao aproximarmos
  $e^x$ pelo seu polinômio de Taylor centrado em zero de grau $n$ é
  \[
    R_n(x) = \frac{\exp^{(n+1)}(\xi)}{(n+1)!} \, x^{n+1}
  \]
  para algum $\xi$ entre zero e $x$.
  Como $\exp^{(k)}(x) = e^x$, para qualquer $k\in\N$, temos que
  \[
    \abs{\exp^{(k)}(t)} \leq e^{\abs{x}}
  \]
  para todo $k\in\N$ e $\abs{t}\leq\abs{x}$.
  Portanto, para todo \(x\in\R\)
  \[
    \abs{R_n(x)}
    = \frac{\abs{\exp^{(n+1)}(\xi)}}{(n+1)!}\abs{x^{n+1}}
    \leq \frac{e^{\abs{x}}}{(n+1)!}\abs{x^{n+1}}
  \]
  Calculando o limite
  \[
    \lim_{n\to\infty} \frac{e^{\abs{x}}}{(n+1)!}\abs{x^{n+1}}
    = e^{\abs{x}}\abs{x^{n+1}} \lim_{n\to\infty} \frac{1}{(n+1)!}
    = e^{\abs{x}}\abs{x^{n+1}} \cdot 0
    = 0
  \]
  podemos usar o \textbf{Teorema do Confronto} para mostrar que \(R_n(x)\)
  vai para zero, quando $n$ vai para o infinito.
  Portanto, a série converge para a função.
\end{answer}

%------------------------------------------------------------------------------%
\end{document}
%------------------------------------------------------------------------------%
